\chapter{The Category Mol}\label{ch:mol}

This chapter builds the algebraic foundation of the thesis. The objects of the
category \textbf{Mol} are types; the morphisms are \emph{stateful
transformations}: deterministic Mealy machines
$(S, \mathrm{step})$ with $\mathrm{step} : S \times A \to B \times S$,
quotiented by the condition that two machines are equal exactly when their
state spaces are in bijection and the bijection carries one transition
structure to the other. State is then an implementation detail: a morphism is
an input--output behavior together with its internal state space, not a
particular machine. On this we establish the category structure (\autoref{thm:1}--\autoref{thm:2}),
the symmetric monoidal structure (\autoref{thm:3}), the subcategories of pure, effectful,
and hybrid molecules (\autoref{thm:4}--\autoref{thm:5}), behavioral equivalence as a congruence (\autoref{thm:6}),
closure under the tensor (\autoref{thm:7}), and determinism (\autoref{thm:8}). Throughout, types are
the types of the ambient universe of \autoref{ch:prelim}, $A \times B$ is the
product type, and $()$ is the unit type.

\section{Machines and Morphisms}\label{sec:mol-machines}

\begin{definition}[Machine]\label{def:machine}
Let $A$ and $B$ be types. A \emph{machine} from $A$ to $B$ is a pair
$(S, \mathrm{step})$ consisting of a type $S$, the \emph{state space}, and a
total function
\[ \mathrm{step} : S \times A \to B \times S . \]
A machine is a deterministic Mealy machine: on input $a$ in state $s$ it
produces output $b$ and next state $s_1$ where
$\mathrm{step} (s, a) = (b, s_1)$ \citep{mealy, moore}. Determinism is built
in: $\mathrm{step}$ is a function, so from a given state and input exactly one
output and one next state result.
\end{definition}

\begin{definition}[State-space isomorphism]\label{def:stiso}
Two machines $(S, \mathrm{step})$ and $(S', \mathrm{step}')$ from $A$ to $B$
are \emph{state-space isomorphic}, written
$(S, \mathrm{step}) \cong (S', \mathrm{step}')$, if there exists a bijection
$\varphi : S \to S'$ such that for all $s \in S$ and $a \in A$,
\[ \mathrm{step} (s, a) = (b, s_1) \;\Longrightarrow\;
   \mathrm{step}'\bigl (\varphi (s), a\bigr) = \bigl (b, \varphi (s_1)\bigr) . \]
\end{definition}

\begin{remark}
State-space isomorphism is an equivalence relation: reflexive with
$\varphi = \mathrm{id}_S$, symmetric with $\varphi^{-1}$, transitive by
composing bijections. Isomorphic machines differ only in the names of their
states: for every input they emit the same outputs and traverse matching
states. This is the \emph{bijection quotient} announced in
\autoref{ch:intro}: state is unobservable in Mol except through its type.
\end{remark}

\begin{definition}[The category Mol]\label{def:mol}
The category \textbf{Mol} has the types as objects. For types $A$ and $B$, a
morphism $f : A \to B$ of Mol is an equivalence class $[S, \mathrm{step}]$ of
machines from $A$ to $B$ under state-space isomorphism. Composition and
identity are defined in \autoref{def:comp} and \autoref{def:id}; that they
are well defined on classes and satisfy the category laws is the content of
\autoref{thm:1} and \autoref{thm:2}.
\end{definition}

\begin{definition}[Composition]\label{def:comp}
Let $f = [S, \mathrm{step}] : A \to B$ and $g = [T, \mathrm{tstep}] : B \to C$
be morphisms of Mol. Their \emph{composite} is the class
\[ g \circ f = [\, T \times S,\; \mathrm{comp} \,] : A \to C , \]
where $\mathrm{comp} : (T \times S) \times A \to C \times (T \times S)$ is given
by
\[ \mathrm{comp}\bigl ((t, s), a\bigr) = \bigl (c, (t_1, s_1)\bigr)
   \quad\text{whenever}\quad (b, s_1) = \mathrm{step} (s, a),\; (c, t_1)
   = \mathrm{tstep} (t, b) . \]
The composite machine runs $f$ on the input, feeds $f$'s output and the
$T$-component of the state to $g$, and records both next states; the state
space of a composite is the product of the state spaces. That the class
$g \circ f$ is independent of the chosen representatives is \autoref{thm:2}.
\end{definition}

\begin{definition}[Identity]\label{def:id}
For each type $A$, the \emph{identity} morphism is
\[ \mathrm{id}_A = [\, (),\; \mathrm{idstep} \,] : A \to A ,
   \qquad \mathrm{idstep}\bigl ((), a\bigr) = (a, ()) . \]
The identity machine has the unit type as state space and copies its input to
the output.
\end{definition}

\begin{theorem}[Mol is a category (\autoref{thm:1})]\label{thm:1}
\textbf{[N].} Composition of $(S, \mathrm{step})$ classes is associative with
identity: Mol is a category.
\end{theorem}

\begin{proof}
Composition on classes is well defined by \autoref{thm:2}, proved below; given
representatives, the composite class is computed as in \autoref{def:comp}. We
verify the category laws on representatives; the quotient inherits them.

\emph{Associativity.} Let $f : A \to B$, $g : B \to C$, $h : C \to D$ be
realized by $(S, \mathrm{step})$, $(T, \mathrm{tstep})$,
$(R, \mathrm{rstep})$. By \autoref{def:comp}, $(h \circ g) \circ f$ and
$h \circ (g \circ f)$ are realized by the machines with state spaces
$(R \times T) \times S$ and $R \times (T \times S)$. Unwinding the definition,
both machines process an input $a$ by running $f$, then $g$, then $h$ in that
order, emitting $h$'s output and recording all three next states; the two
machines differ only in how the state triple is bracketed. Formally, the
canonical bijection
\[ \varphi : (R \times T) \times S \to R \times (T \times S),
   \qquad \varphi\bigl ((r, t), s\bigr) = \bigl (r, (t, s)\bigr) , \]
is a state-space isomorphism: if the first machine takes
$x = (r, t), s)$ on input $a$ to $(d, x_1) = (d, (r_1, t_1), s_1))$, then
the second takes $\varphi (x) = (r, (t, s))$ on $a$ to
$(d, (r_1, (t_1, s_1))) = (d, \varphi (x_1))$, because both machines run the
same three components in the same order on the same data. Hence
$(h \circ g) \circ f = h \circ (g \circ f)$ in Mol.

\emph{Identity laws.} For $f = [S, \mathrm{step}] : A \to B$, the composite
$\mathrm{id}_B \circ f$ is realized by $(() \times S, \mathrm{comp})$, which
runs $f$ and then copies the output, taking $((), s)$ on input $a$ to
$(b, (), s_1))$ when $\mathrm{step} (s, a) = (b, s_1)$. The bijection
$\varphi ((), s) = s$ is a state-space isomorphism onto $f$, since
$\mathrm{step} (\varphi ((), s), a) = \mathrm{step} (s, a) = (b, s_1) =
(b, \varphi ((), s_1))$; so $\mathrm{id}_B \circ f = f$. Symmetrically,
$f \circ \mathrm{id}_A = f$ via $\varphi (s, ()) = s$. Both laws therefore
hold on classes.
\end{proof}

\begin{theorem}[The bijection quotient is well defined (\autoref{thm:2})]\label{thm:2}
\textbf{[N].} The state-space bijection quotient is well defined: equivalent
machines compose to equivalent machines. If
$(S, \mathrm{step}) \cong (S_1, \mathrm{step}_1)$ and
$(T, \mathrm{tstep}) \cong (T_1, \mathrm{tstep}_1)$, then
\[ (T, \mathrm{tstep}) \circ (S, \mathrm{step})
   \;\cong\; (T_1, \mathrm{tstep}_1) \circ (S_1, \mathrm{step}_1) , \]
so the composite class $[g \circ f]$ is independent of the representatives.
\end{theorem}

\begin{proof}
Let $\varphi : S \to S_1$ and $\psi : T \to T_1$ be the state-space
isomorphisms, so that
$\mathrm{step}_1 (\varphi (s), a) = (b, \varphi (s_1))$ whenever
$\mathrm{step} (s, a) = (b, s_1)$, and likewise for $\psi$. The two composites
are realized by $(T \times S, \mathrm{comp})$ and
$(T_1 \times S_1, \mathrm{comp}_1)$ with the transition law of
\autoref{def:comp}. Define $\Phi : T \times S \to T_1 \times S_1$ by
$\Phi (t, s) = (\psi (t), \varphi (s))$; $\Phi$ is a bijection. Suppose
$\mathrm{comp} ((t, s), a) = (c, (t_1, s_1))$, that is,
$(b, s_1) = \mathrm{step} (s, a)$ and $(c, t_1) = \mathrm{tstep} (t, b)$. Then
\[ \mathrm{comp}_1\bigl (\Phi (t, s), a\bigr)
   = \mathrm{comp}_1\bigl (\psi (t), \varphi (s), a\bigr)
   = \bigl (c, (\psi (t_1), \varphi (s_1))\bigr)
   = \bigl (c, \Phi (t_1, s_1)\bigr) , \]
where the middle equality applies the isomorphism hypotheses to the two
component steps. Hence $\Phi$ is a state-space isomorphism of the composites,
which is the claim.
\end{proof}

\begin{remark}
With \autoref{thm:2} in place, composition is a genuine function on the
hom-sets of Mol, and \autoref{thm:1} makes Mol a category. The same
product-of-isomorphisms argument shows that the tensor product of the next
section is well defined on classes.
\end{remark}

\section{Symmetric Monoidal Structure}\label{sec:mol-mon}

The dataplane composes pipelines in two ways: sequentially, by $\circ$, and in
parallel, by a tensor product that pairs independent subpipelines. For Mol the
tensor is the tuple product and the unit is $()$.

\begin{definition}[Tensor product]\label{def:tensor}
On objects, $A \otimes B = A \times B$, and the unit object is
$I = ()$. On morphisms, for $f = [S, \mathrm{step}] : A \to B$ and
$g = [T, \mathrm{tstep}] : C \to D$,
\[ f \otimes g = [\, S \times T,\; \mathrm{step} \otimes \mathrm{tstep} \,]
   : A \times C \to B \times D , \]
where
\[ (\mathrm{step} \otimes \mathrm{tstep})\bigl ((s, t), (a, c)\bigr)
   = \bigl ((b, d), (s_1, t_1)\bigr)
   \quad\text{whenever}\quad (b, s_1) = \mathrm{step} (s, a),\; (d, t_1)
   = \mathrm{tstep} (t, c) . \]
The two machines run independently, one on each component of the input, and
the state space of a tensor is the product of the state spaces.
\end{definition}

\begin{definition}[Braiding]\label{def:braid}
For types $A$ and $B$, the \emph{braiding} is the pure morphism
\[ \sigma_{A,B} = [\, (),\; \mathrm{swap} \,] : A \times B \to B \times A ,
   \qquad \mathrm{swap}\bigl ((), (a, b)\bigr) = \bigl ((b, a), ()\bigr) . \]
\end{definition}

\begin{theorem}[Mol is symmetric monoidal (\autoref{thm:3})]\label{thm:3}
\textbf{[C].} Mol is symmetric monoidal: the tensor product of
\autoref{def:tensor} with unit $()$ is a bifunctor; the associator, the
unitors, and the braiding of \autoref{def:braid} are the canonical pure
morphisms; and they satisfy the coherence conditions of Mac Lane
\citep{maclane}.
\end{theorem}

\begin{proof}
\emph{Well definedness.} If $\varphi : S \to S_1$ and $\psi : T \to T_1$ are
state-space isomorphisms of $f$ with $f_1$ and of $g$ with $g_1$, then
$\Phi (s, t) = (\varphi (s), \psi (t))$ is a state-space isomorphism of
$f \otimes g$ with $f_1 \otimes g_1$: the square commutes componentwise by the
two isomorphism hypotheses, exactly as in the proof of \autoref{thm:2}. The
tensor therefore descends to classes.

\emph{Bifunctoriality.} The tensor of identity machines
$((), \mathrm{idstep}) \otimes (), \mathrm{idstep})$ is realized by
$() \times ()$ with the copy transition, the class
$\mathrm{id}_{A \times B}$. It remains to check the interchange law: for
$f : A \to B$, $f' : B \to C$, $g : A' \to B'$, $g' : B' \to C'$ with
realizations $(S, \mathrm{step})$, $(S', \mathrm{step}')$, $(T, \mathrm{tstep})$,
$(T', \mathrm{tstep}')$,
\[ (f \otimes g) \circ (f' \otimes g')
   = (f \circ f') \otimes (g \circ g') . \]
The left-hand side is realized by the state space
$(S' \times T') \times (S \times T)$: it runs $f'$ and $g'$ on the input pair,
then $f$ and $g$ on the outputs. The right-hand side is realized by
$(S \times S') \times (T \times T')$: it runs $f'$ then $f$, and $g'$ then
$g$. Both machines apply the four transitions $f'$, $g'$, $f$, $g$ exactly
once each, in that order, to the corresponding components, and the canonical
bijection
\[ \Phi\bigl ((s', t'), (s, t)\bigr) = \bigl ((s, s'), (t, t')\bigr) \]
carries the state after the four transitions of the first machine to the
state after the four transitions of the second; the outputs coincide
componentwise. Hence $\Phi$ is a state-space isomorphism and the interchange
law holds, so $\otimes$ is a bifunctor \citep{maclane}.

\emph{Naturality of the braiding.} For $f : A \to A'$ and $g : B \to B'$
with realizations $(S, \mathrm{step})$, $(T, \mathrm{tstep})$, both
$(g \otimes f) \circ \sigma_{A,B}$ and
$\sigma_{A',B'} \circ (f \otimes g)$ run $f$ on the first component and $g$
on the second, then swap the outputs; their state spaces are
$T \times S$ and $S \times T$, matched by the bijection
$\Phi (t, s) = (s, t)$, which commutes with the transitions because the swap
step is pure. Hence $\sigma$ is natural in both arguments.

\emph{Associator and unitors.} The associator
$\alpha_{A,B,C} : (A \times B) \times C \to A \times (B \times C)$ is the pure
machine that reassociates the tuple,
$\alpha ((), (a, b), c)) = (a, (b, c)), ())$; the left and right unitors
$\lambda_A : () \times A \to A$ and $\rho_A : A \times () \to A$ are the pure
projection machines, $\lambda ((), (), a)) = (a, ())$ and
$\rho ((), (a, ())) = (a, ())$. Each is an isomorphism in Mol, its inverse
being the pure machine running the same tuple manipulation backwards.

\emph{Coherence.} The pentagon, triangle, and hexagon diagrams commute
because both paths through each diagram compute the same function on tuples.
In the pentagon, both composite paths reassociate
$((A \otimes B) \otimes C) \otimes D$ to
$A \otimes (B \otimes (C \otimes D))$, and reassociation of tuples is unique.
In the triangle, both paths project $((A \otimes I) \otimes B)$ to
$A \otimes B$; in the hexagon, both sides of
\[ (1_B \otimes \sigma_{A,C}) \circ (\sigma_{A,B} \otimes 1_C)
   = \sigma_{A,B \otimes C} \]
map $(a, b, c)$ to $(b, c, a)$, since the left-hand side first swaps $a$ and
$b$ and then swaps $a$ and $c$, and the right-hand side swaps the single
$A$-component past the pair $B \otimes C$. Similarly
$\sigma_{B,A} \circ \sigma_{A,B} = \mathrm{id}_{A \otimes B}$, since two
swaps return the tuple. Hence $\sigma$ is a coherent symmetry. By Mac Lane's
coherence theorem every diagram whose edges are built from $\alpha$, $\lambda$,
$\rho$, and $\sigma$ commutes, so Mol is symmetric monoidal: and, since
composition is by \autoref{thm:2}, the coherence equations hold on classes
\citep{maclane}. This is the structure on which the free constructions of
\autoref{ch:free} and the traced structure of \autoref{ch:trace} are built.
\end{proof}

\section{Pure Molecules}\label{sec:mol-pure}

A molecule with no state at all computes a total function. The subcategory
these form is the functional fragment of Mol.

\begin{definition}[PureMol]\label{def:puremol}
\textbf{PureMol} is the subcategory of Mol with the same objects and with
morphisms the classes of machines representable with state space $()$: a
morphism $f : A \to B$ of PureMol is a class $[(), \mathrm{step}]$. The
quotient is trivial on pure machines: because the only bijection
$() \to ()$ is the identity, two pure classes are equal exactly when their
step functions agree pointwise.
\end{definition}

\begin{theorem}[PureMol is the category of functions (\autoref{thm:4})]\label{thm:4}
\textbf{[N].} PureMol is equivalent to the category \textbf{Type} of types
and total functions between them; in fact the two categories are isomorphic.
\end{theorem}

\begin{proof}
Define $E : \mathrm{PureMol} \to \mathbf{Type}$ as the identity on objects
and, on morphisms,
$E[(), \mathrm{step}] = f$ where $\mathrm{step} ((), a) = (f (a), ())$. This is
well defined by \autoref{def:stiso}: two pure machines are state-space
isomorphic exactly when their step functions agree pointwise. It is a
functor: by \autoref{def:comp} the composite of the pure machines for $f$ and
$g$ runs $f$ then $g$ and hence represents $g \circ f$, and the identity
machine represents $\mathrm{id}_A$. $E$ is full: every total function
$f : A \to B$ is realized by the machine $((), \mathrm{step}_f)$ with
$\mathrm{step}_f ((), a) = (f (a), ())$. $E$ is faithful: if
$E[(), \mathrm{step}] = E[(), \mathrm{step}']$, the step functions agree, so
the classes are equal. And $E$ is bijective on objects. Hence $E$ is an
equivalence of categories \citep{maclane}; since the inverse on hom-sets
sends $f$ to $[(), \mathrm{step}_f]$ and is the identity on objects, the
equivalence is an isomorphism of categories. PureMol is the type
theoretic shadow of Mol itself: objects as types, morphisms as functions.
\end{proof}

\section{Effectful Molecules and the Kleisli Category}\label{sec:mol-eff}

A dataplane carries a fixed ambient \emph{context} $\mathrm{Ctx}$: the
packet, the workspace, the request: through every pipeline stage. An
effectful molecule is a machine whose only state is the context; composing
effectful molecules threads a single context through both, which is exactly
the bind of the state monad.

\begin{definition}[EffMol]\label{def:effmol}
Fix a type $\mathrm{Ctx}$ of contexts. \textbf{EffMol}$(\mathrm{Ctx})$ is the
category with the same objects as Mol whose morphisms $A \to B$ are the step
functions $\mathrm{step} : A \times \mathrm{Ctx} \to B \times \mathrm{Ctx}$
These are machines with state space $\mathrm{Ctx}$: compared by literal equality
of functions. The identity on $A$ is $\mathrm{id} (a, s) = (a, s)$, and the
composite of $\mathrm{step} : A \times \mathrm{Ctx} \to B \times \mathrm{Ctx}$
with $\mathrm{tstep} : B \times \mathrm{Ctx} \to C \times \mathrm{Ctx}$ is the
\emph{context-threading} function
\[ (\mathrm{tstep} \circ \mathrm{step})(a, s) = (c, s_2)
   \quad\text{where}\quad (b, s_1) = \mathrm{step} (a, s),\; (c, s_2)
   = \mathrm{tstep} (b, s_1) . \]
Both machines see the same single context, threaded from first to second.
\end{definition}

\begin{remark}
Every step function of EffMol$(\mathrm{Ctx})$ is a machine with state space
$\mathrm{Ctx}$ and hence determines a class $[\mathrm{Ctx}, \mathrm{step}]$
of Mol; but EffMol is not a subcategory of Mol under inherited composition.
The Mol-composite of two context machines has state space
$\mathrm{Ctx} \times \mathrm{Ctx}$(the contexts multiply) while the
EffMol-composite threads one copy; and the EffMol identity is the class of a
machine that is the Mol identity only when $\mathrm{Ctx}$ is a singleton.
EffMol is the model of the effectful fragment; its relation to Mol is
developed in \autoref{ch:kleisli}.
\end{remark}

\begin{theorem}[EffMol is the Kleisli category of the state monad (\autoref{thm:5})]\label{thm:5}
\textbf{[N].} EffMol$(\mathrm{Ctx})$ is isomorphic to the Kleisli category of
the state monad on $\mathrm{Ctx}$.
\end{theorem}

\begin{proof}
Recall the state monad of \autoref{ch:prelim}: on the category of types,
$T X = (X \times \mathrm{Ctx})^{\mathrm{Ctx}}$ with unit
$\eta_X (x) = \lambda s.\,(x, s)$ and multiplication threading two state
computations. Both EffMol$(\mathrm{Ctx})$ and the Kleisli category
$\mathrm{Kl} (T)$ have the types as objects. A Kleisli morphism $A \to B$ is a
function $A \to T B = (B \times \mathrm{Ctx})^{\mathrm{Ctx}}$, which by
currying is exactly a function $A \times \mathrm{Ctx} \to B \times
\mathrm{Ctx}$, i.e.\ a morphism of EffMol$(\mathrm{Ctx})$; this bijection
$F$ on hom-sets is the identity on objects. The Kleisli identity is
$\eta_A$, which under currying is the EffMol identity $(a, s) \mapsto (a, s)$.
Kleisli composition is the bind $g \circ_T f = \mu_C \circ T g \circ f$, which
uncurries to the context-threading law
\[ (g \circ_T f)(a, s) = (c, s_2) \quad\text{where}\quad (b, s_1) = f (a, s),\; (c, s_2) = g (b, s_1) , \]
matching \autoref{def:effmol} exactly. Hence $F$ preserves identity and
composition; the inverse construction sends each step function to its
curried form. $F$ is therefore an isomorphism of categories
\citep{maclane, arxiv-monads-kleisli}. The parenthetical ``algebras over the
monad'' of the catalog refers to the Eilenberg--Moore category, the other
resolution of the same monad; the isomorphism here is with the Kleisli
category, which is what the dataplane threads.
\end{proof}

\section{Hybrid Molecules}\label{sec:mol-hybrid}

Real dataplane stages carry both an explicit pure state: a counter, a
checksum register, a ring pointer: and the ambient context. A hybrid
molecule is a machine whose state space splits as a product of a pure part
and the context.

\begin{definition}[HybridMol]\label{def:hybridmol}
Fix the context type $\mathrm{Ctx}$. \textbf{HybridMol}$(\mathrm{Ctx})$ is
the subcategory of Mol with the same objects and with morphisms the classes
of machines $(S, \mathrm{step})$ whose state space $S$ is a singleton or is
isomorphic to $S_{\mathrm{pure}} \times \mathrm{Ctx}$ for some type
$S_{\mathrm{pure}}$. The singleton clause includes the pure identity and with
it all pure morphisms, which the identity laws of \autoref{thm:1} require.
Equivalently, a hybrid molecule is a machine with state space
$S_{\mathrm{pure}} \times \mathrm{Ctx}$: a pure state component
$S_{\mathrm{pure}}$ that the stage owns and a context component
$\mathrm{Ctx}$ that it shares with the ambient pipeline.
\end{definition}

\begin{remark}
HybridMol is closed under composition and tensor. The composite of hybrid
machines with pure parts $S_{\mathrm{pure}}$ and $T_{\mathrm{pure}}$ is
realized by
\[ (T_{\mathrm{pure}} \times \mathrm{Ctx}) \times (S_{\mathrm{pure}} \times \mathrm{Ctx})
   \;\cong\; (T_{\mathrm{pure}} \times S_{\mathrm{pure}} \times \mathrm{Ctx})
   \times \mathrm{Ctx} , \]
again of the form (pure part) $\times$ (context), the canonical isomorphism
being a state-space isomorphism by the product case of \autoref{thm:2}; the
tensor closes by the same computation, absorbing the extra context copy into
the pure part, as recorded in \autoref{thm:7}. EffMol$(\mathrm{Ctx})$ sits
inside HybridMol as the case $S_{\mathrm{pure}} = ()$. The decomposition of a
hybrid molecule as $M = q \circ e \circ p$ with $p, q$ pure and $e$
effectful is proved in \autoref{ch:kleisli}.
\end{remark}

\section{Behavioral Equivalence}\label{sec:mol-bisim}

Two machines may differ in their state spaces yet compute the same
transformation on input words. Behavioral equivalence identifies exactly
those machines, and it is the relation under which Mol is quotiented in the
minimization and testing chapters.

\begin{definition}[Response]\label{def:response}
For a machine $M = (S, \mathrm{step}) : A \to B$ and a state $s \in S$, the
\emph{response} of $s$ is the function $r_{M,s} : A^{*} \to B^{*}$ defined by
$r_{M,s} (\varepsilon) = \varepsilon$ and, for an input symbol $a$ and a word
$w$ over $A$,
\[ r_{M,s} (a \cdot w) = b \cdot r_{M,s_1} (w)
   \quad\text{wherever}\quad \mathrm{step} (s, a) = (b, s_1) , \]
with $\cdot$ word concatenation. The response is the full input--output
behavior of the machine started in state $s$.
\end{definition}

\begin{definition}[Behavioral equivalence]\label{def:bisim}
Machines $M, M' : A \to B$ are \emph{behaviorally equivalent}, written
$M \approx M'$, if their sets of realized responses coincide:
\[ \{ r_{M,s} : s \in S \} = \{ r_{M',s'} : s' \in S' \} . \]
Each state of one machine realizes a response realized by some state of the
other. For deterministic machines this relation coincides with bisimulation
in the sense of Park \citep{park}.
\end{definition}

\begin{remark}
State-space isomorphic machines are behaviorally equivalent: the isomorphism
pairs states with equal responses, and the response of the image state is
the response of the source state by \autoref{def:stiso}. Hence $\approx$ is
well defined on the morphisms of Mol. It is an equivalence relation, being
equality of sets, and it is strictly coarser than state-space isomorphism: a
machine with an unused state realizes the same responses as the machine
without it.
\end{remark}

\begin{theorem}[Behavioral equivalence is a congruence (\autoref{thm:6})]\label{thm:6}
\textbf{[N].} Behavioral equivalence (bisimulation) is a congruence for
$\circ$ and $\otimes$; the quotient category $\mathrm{Mol}/{\approx}$ is a
symmetric monoidal category.
\end{theorem}

\begin{proof}
\emph{Composition is a congruence.} First observe that the response of a
state determines the responses of all its successors: if
$r_{M,s} = r_{M,s'}$, $\mathrm{step} (s, a) = (b, s_1)$, and
$\mathrm{step} (s', a) = (b, s_1')$, then for every word $w$,
$b \cdot r_{M,s_1} (w) = r_{M,s} (a \cdot w) = r_{M,s'} (a \cdot w) =
b \cdot r_{M,s_1'} (w)$, so $r_{M,s_1} = r_{M,s_1'}$. Now let $f, g$ be
realized by $(S, \mathrm{step})$, $(T, \mathrm{tstep})$ and write
$R_f = \{ r_{f,s} \}$, $R_g = \{ r_{g,t} \}$. Reading an input word symbol by
symbol, the composite machine runs $f$ and threads its outputs into $g$; by
the observation, the response of the composite at $(t, s)$ is determined by
the pair $(r_{g,t}, r_{f,s})$ alone, through the recurrence
\[ r_{g \circ f, (t,s)} (a \cdot w)
   = c \cdot r_{g \circ f, (t_1, s_1)} (w) , \]
where $(b, s_1) = \mathrm{step} (s, a)$, $(c, t_1) = \mathrm{tstep} (t, b)$,
and the successor responses in the recursion are functions of
$(r_{g,t}, r_{f,s})$ by the observation. Hence the response set of $g \circ f$
depends only on $R_f$ and $R_g$; if $f \approx f'$ and $g \approx g'$ then
$R_f = R_{f'}$ and $R_g = R_{g'}$, so the two composites have the same
response set: $g \circ f \approx g' \circ f'$.

\emph{Tensor is a congruence.} The tensor machine runs $f$ and $g$
independently on the two components, so
\[ r_{f \otimes g, (s,t)} (w_1, w_2) = \bigl (r_{f,s} (w_1),\, r_{g,t} (w_2)\bigr) , \]
and the response set of $f \otimes g$ is $R_f \times R_g$. Equality of the
component sets gives $f \otimes g \approx f' \otimes g'$. Hence $\approx$ is
a congruence for both operations.

\emph{The quotient.} A congruence is exactly what is needed for the
operations to descend to classes: define
$[g]_{\approx} \circ [f]_{\approx} = [g \circ f]_{\approx}$ and
$[f]_{\approx} \otimes [g]_{\approx} = [f \otimes g]_{\approx}$, well defined
by the two congruences. The category laws and the coherence equations of
\autoref{thm:3} hold in Mol and are equations between morphisms, so they hold
in the quotient; the structural morphisms of \autoref{thm:3} remain
isomorphisms there, since e.g.
$[\sigma_{B,A}]_{\approx} \circ [\sigma_{A,B}]_{\approx} =
[\mathrm{id}]_{\approx}$. Hence $\mathrm{Mol}/{\approx}$ is a symmetric
monoidal category. On small finite machines the congruence property is
checked exhaustively by the \texttt{bisim} tool of \autoref{app:a}.
\end{proof}

\begin{theorem}[Closure under the tensor (\autoref{thm:7})]\label{thm:7}
\textbf{[N].} PureMol and EffMol are closed under the tensor: the tensor of
two pure molecules is pure; the tensor of an effectful molecule of
EffMol$(\mathrm{Ctx})$ with a pure molecule, on either side, is effectful;
and the tensor of two effectful molecules of EffMol$(\mathrm{Ctx})$ is a
molecule of EffMol$(\mathrm{Ctx} \times \mathrm{Ctx})$, equivalently a hybrid
molecule.
\end{theorem}

\begin{proof}
The tensor of pure machines $((), \mathrm{step})$ and
$((), \mathrm{tstep})$ is realized by the state space $() \times ()$ with the
product transition; the canonical bijection $() \times () \to ()$ is a
state-space isomorphism, so the tensor class is pure. For
$f = [\mathrm{Ctx}, \mathrm{step}]$ in EffMol$(\mathrm{Ctx})$ and a pure
machine $((), \mathrm{tstep})$, the tensor $f \otimes p$ is realized by
$\mathrm{Ctx} \times ()$; the bijection $\Phi (s, ()) = s$ is a state-space
isomorphism onto the machine with state space $\mathrm{Ctx}$ and transition
$\mathrm{step}_1 (s, (a, c)) = (b, d), s_1)$ where $(b, s_1) =
\mathrm{step} (s, a)$ and $d$ is the output of the pure component on $c$,
because the pure component performs no state update. Hence
$f \otimes p \in \mathrm{EffMol} (\mathrm{Ctx})$; the case $p \otimes f$ is
symmetric. Finally, the tensor of two context machines is realized by
$\mathrm{Ctx} \times \mathrm{Ctx}$, which is a morphism of
EffMol$(\mathrm{Ctx} \times \mathrm{Ctx})$ and, being of the form
$S_{\mathrm{pure}} \times \mathrm{Ctx}$ with $S_{\mathrm{pure}} =
\mathrm{Ctx}$, a hybrid molecule by \autoref{def:hybridmol}.
\end{proof}

\begin{remark}
The exact boundary of \autoref{thm:7} is worth stating: the tensor of two
EffMol$(\mathrm{Ctx})$ molecules has state space $\mathrm{Ctx} \times
\mathrm{Ctx}$, and since $\mathrm{Ctx} \times \mathrm{Ctx} \cong \mathrm{Ctx}$
fails in general (for finite $\mathrm{Ctx}$ with more than one element),
EffMol$(\mathrm{Ctx})$ is not closed under the tensor with itself in the
strict sense. What the tensor always preserves is purity strictly and
effectfulness in the graded sense above: contexts add, they do not vanish.
\end{remark}

\section{Atoms and Determinism}\label{sec:mol-atoms}

Molecules are built from \emph{atoms}: the primitive components of a
dataplane, total on their declared domain, with no hidden state beyond their
declared one.

\begin{definition}[Atoms]\label{def:atoms}
An \emph{atom} is a molecule that is a generator rather than a composite:
formally, a chosen morphism of a signature $\Sigma$, and a \emph{molecule
over} $\Sigma$ is any morphism of the subcategory of Mol generated by
$\Sigma$ under $\circ$ and $\otimes$ \citep{lawvere}. Concretely, the atoms
of a dataplane are the pure atoms: total functions such as \emph{parse},
\emph{checksum}, \emph{push}, \emph{pop}: and the effectful atoms of
EffMol$(\mathrm{Ctx})$, such as reading and writing context fields. Every
atom is total on its declared domain and deterministic. The theories of
signatures are the subject of \autoref{ch:lawvere}, and their free molecules
and normal forms of \autoref{ch:free}.
\end{definition}

\begin{theorem}[Determinism (\autoref{thm:8})]\label{thm:8}
\textbf{[N].} Pure molecules and tensor products of pure molecules are
deterministic; behavioral equivalence preserves determinism.
\end{theorem}

\begin{proof}
A molecule is deterministic by construction of \autoref{def:machine}: its
transition function is a total function, so from any state and input exactly
one output and one next state result; no Mol morphism involves choice. In
particular every pure molecule is deterministic, and the tensor of pure
molecules is pure by \autoref{thm:7}, hence deterministic. The content of the
theorem is that determinism, unlike purity, is a property of behavior: for
every machine $M$ and state $s$, the response $r_{M,s}$ of \autoref{def:response}
is a total single-valued function $A^{*} \to B^{*}$. Behavioral equivalence
compares exactly these functions (\autoref{def:bisim}), so if
$M \approx M'$ then $M'$ realizes precisely the responses of $M$: a
deterministic machine has only deterministic equivalents, and every class of
$\mathrm{Mol}/{\approx}$ consists of deterministic behaviors. Purity, by
contrast, is not invariant under $\approx$: a machine with an unused state is
behaviorally equivalent to a pure machine without it, yet is not pure. This
is why the rewrite and minimization chapters normalize behavior first
(\autoref{ch:rewrite}, \autoref{ch:coalgebra}): determinism is preserved by
every equivalence the engine uses, and the same input always yields the same
output through the same states.
\end{proof}

\section{Discussion}\label{sec:mol-concl}

Mol realizes the design of \autoref{ch:intro}: a category whose objects are
types and whose morphisms are stateful transformations, quotiented so that
state is an implementation detail, and closed under the two composition
schemes of the dataplane. The category laws (\autoref{thm:1}--\autoref{thm:2}) license equational
reasoning about pipelines; the symmetric monoidal structure (\autoref{thm:3}) licenses
parallel composition and the tuple manipulations used throughout
\autoref{ch:lawvere}--\autoref{ch:rewrite}; the subcategories PureMol,
EffMol$(\mathrm{Ctx})$, and HybridMol$(\mathrm{Ctx})$
(\autoref{sec:mol-pure}--\autoref{sec:mol-hybrid}) fix the vocabulary in
which \autoref{ch:kleisli} proves the pure/effectful decomposition; and
behavioral equivalence (\autoref{thm:6}) is the relation under which minimization
(\autoref{ch:coalgebra}) and testing (\autoref{ch:testing}) operate, while
determinism (\autoref{thm:8}) is the assumption under which the feedback structure of
\autoref{ch:trace} and the decision problems of \autoref{ch:situations} are
stated. The next chapter adds equational theories: dataplane components are
models of Lawvere theories, and the ring, parser, and transport equations
become laws that molecules must satisfy \citep{lawvere}.

\begin{table}[htbp]
\centering
\small
\begin{tabular}{@{}llp{9.3cm}@{}}
\toprule
Theorem & Tag & Statement \\
\midrule
\ref{thm:1} & [N] & Mol is a category: composition of $(S, \mathrm{step})$ classes is associative with identity. \\
\ref{thm:2} & [N] & The state-space bijection quotient is well defined: equivalent machines compose to equivalent machines. \\
\ref{thm:3} & [C] & Mol is symmetric monoidal: tensor the tuple product, unit $()$, coherent braiding \citep{maclane}. \\
\ref{thm:4} & [N] & PureMol is equivalent to the category of types and total functions. \\
\ref{thm:5} & [N] & EffMol$(\mathrm{Ctx})$ is isomorphic to the Kleisli category of the state monad on $\mathrm{Ctx}$. \\
\ref{thm:6} & [N] & Behavioral equivalence is a congruence for $\circ$ and $\otimes$; Mol$/{\approx}$ is symmetric monoidal. \\
\ref{thm:7} & [N] & PureMol and EffMol are closed under the tensor. \\
\ref{thm:8} & [N] & Pure molecules and tensor products of pure molecules are deterministic; behavioral equivalence preserves determinism. \\
\bottomrule
\end{tabular}
\caption{The theorems of this chapter; the statements reproduce the catalog of \autoref{sec:contributions}. [N] marks statements original to this thesis, proved with elementary means; [C] marks classical results restated and fully proved inside Mol, with the source cited.}
\end{table}

\section{Optimality of the model}\label{sec:mol-optimality}

This section makes precise the claim that the pure, effectful, and
hybrid classes of Mol are forced by the model rather than chosen by
taste, and that the model excludes no application. Two theorems
establish the claim: the classification is the unique extension of the
atom classification that respects the algebraic structure
(\autoref{thm:classification}), and the free construction is initial,
so every interpretation factors through it and no application is
rigidified (\autoref{thm:universality}).

\begin{theorem}[The classification is forced]\label{thm:67}\label{thm:classification}
Let $\Sigma = \Sigma_p \uplus \Sigma_e$ be an atom signature partitioned
into pure and effectful atoms, and let $F(\Sigma)$ be the free
symmetric monoidal category on $\Sigma$. There is exactly one
assignment $c : \mathrm{Mor}(F(\Sigma)) \to \{P, E, H\}$ of the three
classes that (i) agrees with the partition on atoms, (ii) assigns the
structural morphisms (identities, braidings) to $P$, and (iii) is a
congruence: the class of a composite or tensor is a function of the
classes of the factors, with the closure rules
$P \circ P \subseteq P$, $P \circ E \subseteq E$,
$E \circ P \subseteq E$, $E \circ E \subseteq E$ (sequential
composition is Kleisli-closed), $P \otimes P \subseteq P$,
$P \otimes E \subseteq E$, $E \otimes E \subseteq H$, and
$H \otimes \cdot \subseteq H$. The class map is
\[
c(f) = \begin{cases}
P & \text{if } f \text{ is pure,}\\
E & \text{if } f \text{ is effectful,}\\
H & \text{otherwise.}
\end{cases}
\]
Every class of the free model is one of the three; no fourth class is
generated, and no two are identified.
\end{theorem}

\begin{proof}
Define $c$ by induction on the syntax of morphisms of $F(\Sigma)$.
Atoms: $c(a) = P$ for $a \in \Sigma_p$ and $c(a) = E$ for
$a \in \Sigma_e$; structural morphisms have $c = P$. For composites,
apply the closure rules; for tensors, apply the tensor closure rules.
The assignment is total by induction, because every morphism of the
free category is built from atoms, structural morphisms, $\circ$, and
$\otimes$. It is well-defined with respect to the equational theory of
$F(\Sigma)$: the symmetric monoidal axioms (associativity, unit,
interchange, braiding) equate morphisms whose class assignment agrees,
because the class of a composite is a function of the classes of its
factors and the closure rules are closed under the structural
equations (identities and braidings are pure, and composing or
tensoring with a pure morphism preserves the class by the closure
rules). Uniqueness: any assignment satisfying (i)--(iii) must agree
with $c$ on atoms and structural morphisms, and must then agree on
every composite by induction on its syntax, because (iii) forces the
class of each composite to be the closure-rule value of the classes of
its factors. The ``no fourth class'' claim is the totality of the
three-valued map; the ``no identification'' claim follows because the
map is well-defined, so two morphisms in different classes cannot be
equated by the theory.
\end{proof}

\begin{theorem}[Universality: no application is excluded]\label{thm:68}\label{thm:universality}
(1) For any symmetric monoidal category $\mathcal{C}$ and any
interpretation of the signature $\Sigma$ in $\mathcal{C}$, there is a
unique strict symmetric monoidal functor
$F(\Sigma) \to \mathcal{C}$ extending the interpretation.
(2) Every finite-state Mealy transformation is realized by a molecule
over a finite signature. (3) Mol is closed under products, sums,
feedback (traces), and batching. Consequently every application
expressible as a deterministic stateful transformation is expressible
as a molecule, and structural equality in the free model is decided by
the coherence theorem of Mac Lane \citep{maclane}.
\end{theorem}

\begin{proof}
(1) is the universal property of the free symmetric monoidal category:
$F(\Sigma)$ is the initial object in the category of symmetric
monoidal categories equipped with a $\Sigma$-indexed family of
objects, and the unique functor is given by interpreting generators as
generators, composites as composites, and coherence isomorphisms by
the coherence theorem. (2) Let $(S, \mathit{step})$ with
$\mathit{step} : S \times A \to B \times S$ be a finite-state Mealy
machine. Take the signature $\Sigma$ with one effectful atom
$e_s : A \to B$ for each state $s \in S$ whose step function is
$e_s(a) = (b, s')$ where $\mathit{step}(s, a) = (b, s')$, and one
choice atom for the initial state. The molecule $\bigotimes$-wiring of
the atoms together with the trace construction of
\autoref{ch:trace} realizes the machine: the state is threaded through
the trace, and each step is the effectful atom for the current state.
The realization is exact because Mealy determinism makes the output a
function of the state and input alone. (3) Products and sums are
objects of the free category over a signature extended by the
corresponding constructors; feedback is the traced structure of
\autoref{ch:trace}; batching is the operadic structure of
\autoref{ch:operad}. The closure statements are theorems of those
chapters.
\end{proof}

\begin{remark}[On the choice of language]
The model uses categories of machines rather than the alternative
frameworks for effectful computation. Each alternative is a proper
substructure of the same picture: the Kleisli category of the state
monad is exactly the effectful subcategory of Mol
(\autoref{sec:mol-eff}); Freyd categories and Hughes arrows are
abstract machines with a premonoidal tensor, while Mol's tensor is a
true product, which is what licenses the interchange and coherence
laws of \autoref{sec:mol-mon}; coalgebraic final semantics underlies
the minimization theorems of \autoref{ch:coalgebra}; effectus theory
targets probabilistic and quantum effects and is not needed for
deterministic stateful transformations; and process calculi
introduce nondeterminism, which would invalidate the cost and
allocation theorems of \autoref{ch:enrich} and \autoref{ch:linear}.
The trichotomy of \autoref{thm:classification} is finer than the
deterministic/probabilistic dichotomy of effectus theory within the
deterministic world, and \autoref{thm:universality} shows that no
application is excluded. The choice of Mol is therefore not a
heuristic preference; it is the minimal language in which the
engineering properties of this thesis are theorems.
\end{remark}
